\section{Delta-Based Workspaces}
\label{sec:delta}

\subsection{Shared baseline, private state}

Delta-based workspaces separate a task's modifications from the state it shares
with other tasks. A workspace binds a shared immutable lower $S$ and maintains
a private writable upper $\Delta_i$; its view is
$W_i=\operatorname{Resolve}(S,\Delta_i)$. The lower may be a sealed catalog in
the metadata database or a packed snapshot on object storage. These choices
change how inherited records are retrieved, while the workspace contract
determines which records are visible. Neither Packed Metadata nor a particular
object-storage layout is required by this contract.

Creating a workspace registers its base identity and independent private
state. The catalog path identifies the base by its layer, sealed version,
and root hash; a packed lower additionally binds its authenticated manifest.
Unchanged entries and file contents remain reachable through this identity,
so creation need not copy the complete metadata tree. Binding validation,
ownership records, and mount preparation still have costs.

A private name record distinguishes a present entry, an explicit deletion,
and absence. For parent directory $d$ and name $n$,
\begin{equation}
R_i(d,n)=
\begin{cases}
e, & \Delta_i(d,n)=\operatorname{Present}(e),\\
\bot, & \text{the name is hidden in }W_i,\\
R_S(d,n), & \text{otherwise}.
\end{cases}
\end{equation}
A whiteout suppresses the inherited name; applicable directory-opacity rules
can also prevent inheritance. Absence permits fallback only when no such rule
hides the name. Removing an upper payload therefore does not undo a deletion
marker and resurrect a lower file.

Directory enumeration applies the same precedence to each name, removes
duplicates, and returns attributes of the visible inode. Private attributes
describe its current state; they are not an arithmetic merge with lower
attributes. Hard-link aliases must continue to identify the same logical
inode when private state is materialized. Namespace mutations such as rename
must update related entries consistently within the workspace.

\subsection{Changed data and local dirty reads}

The catalog path records private data extents alongside namespace and inode
changes. The Entity path supplies one concrete current-state representation:
each changed logical chunk has a normalized map with explicit \emph{Data},
\emph{LowerRange}, and \emph{Hole} spans. Data references newly stored upper
bytes, LowerRange references inherited bytes in $S$, and Hole reads as zeros.
An uncovered interval in an explicit map is not an implicit lower reference.
The map must describe the applicable range completely.

A short overwrite stores new bytes and replaces their mapped interval;
unchanged intervals can keep their lower references. This requires neither a
byte-difference algorithm nor a copy of the whole lower file. Truncation needs
an additional inheritance boundary. Entity maintains a lower-visible size
that can decrease but cannot increase. If a truncated file is extended, the
newly exposed region contains zeros or new upper data, rather than old bytes
resurrected from $S$. Object formation and integrity checking remain separate
choices: a short logical write need not produce an equally short object.

The VFS dirty overlay serves a different purpose: reading this client's bytes
before writeback completes. If local dirty data covers $[4,8)$ KiB, a read of
$[2,10)$ KiB combines two committed 2 KiB ranges with the dirty 4 KiB range.
Workspace delta state instead determines persistent namespace, attribute, and
data-mapping visibility. A buffered workspace path can use both mechanisms;
local read-your-writes behavior alone establishes neither remote durability
nor a published recovery point.

\subsection{Publication and workspace lifetime}

Data publication must respect the selected durability boundary. The Entity
path uploads an object durably before publishing its mapping. The ordinary
VFS offers upload-before-commit and explicitly selected commit-before-upload
modes; the latter can expose a mapping before its object is remotely
available. Consequently, preparing a recoverable version requires an explicit
drain and validation boundary, even when local reads already return new data.

A worker handoff starts with application checkpoints and a stop to new source
writes. Required payloads and metadata must become durable and validated
before an immutable version reference is published. The destination then
binds that version and acquires authority for subsequent private writes.
Lease and generation checks must reject stale writers. The handoff protocol
must establish a file recovery point. The application runtime is responsible
for restoring sessions, restarting tools, and reopening files. The file version
does not transfer running processes or guarantee recovery of progress beyond
the last confirmed publication. The available
publication and recovery paths are distinguished in
Section~\ref{sec:implementation}.

Lifetime management applies to either lower representation. A workspace or
reader holding $S$ must retain its referenced metadata and data. Publishing a
replacement must not invalidate existing readers, and reclamation must
separate live references from abandoned construction objects. For one shared
baseline, live logical metadata obeys
\begin{equation}
M_{\mathrm{live\text{-}meta}}=M_S+\sum_i M_{\Delta_i}+M_{\mathrm{control}}.
\end{equation}
$M_S$ counts the shared metadata once; $M_{\Delta_i}$ counts private names,
attributes, deletion records, and mappings; $M_{\mathrm{control}}$ counts
disjoint binding and ownership records. Upper payload, encoding and checksum overhead,
unreclaimed objects, and caches belong to separate physical accounting.
Private state includes complete current records and mapping overhead, so its
cost grows with modification and deletion. The evaluation will measure this cost
against materialized workspaces under matched visibility and durability
contracts.
